% !TeX root = main.tex

调试的本质是“在程序运行时观察它”，而 Visual Studio 的调试器把这件苦差事变成了一套可以自由组合的工具：断点决定在哪里停下，单步决定怎么走，数据窗口决定看什么，调用堆栈决定“我是从哪儿来的”。这一章按实际使用顺序把 C++ 调试的功能串成一条线，从最基础的 \key{F5} 一路讲到数据断点、转储分析与多线程死锁排查。读完它，你应该能对任何一个崩溃或数值异常说清楚“下一步该打开哪个窗口”。

\section{调试的基本流程}

\subsection{调试运行与“不调试运行”}

\key{F5} 是最常用的启动方式：\ui{调试 > 开始调试}。它会先编译，再把调试器附加到新启动的进程上，加载符号，然后停在第一个命中的断点处。此时进程处于\textbf{中断模式}（break mode），所有调试窗口才有内容；按 \key{F5} 继续后回到\textbf{运行模式}，大部分窗口会变灰。会话控制还有三个命令：\key{Shift+F5} 停止调试，\key{Ctrl+Shift+F5} 重启调试（比“停止再 F5”快，因为不必重新附加），\ui{调试 > 全部中断} 用来打断正在死循环或死锁的程序。\key{Ctrl+F5}（\ui{调试 > 开始执行（不调试）}）则只是把程序跑起来，不附加调试器，二者的差别远不止“能不能下断点”：

\begin{longtable}{L{3.0cm}L{5.8cm}L{5.8cm}}\toprule
\textbf{对比项} & \textbf{\key{F5} 调试运行} & \textbf{\key{Ctrl+F5} 不调试运行} \\ \midrule \endfirsthead
\toprule \textbf{对比项} & \textbf{\key{F5} 调试运行} & \textbf{\key{Ctrl+F5} 不调试运行} \\ \midrule \endhead \bottomrule \endlastfoot
调试器 & 已附加，可随时中断 & 未附加 \\
断点 / 单步 / 异常 & 有效，可在异常抛出时中断 & 完全无效，只按程序自身逻辑处理 \\
调试堆与符号 & 启用调试堆，按设置加载 .pdb & 都不启用 \\
速度与适用场景 & 明显更慢，用于定位问题 & 接近真实速度，用于只看结果 \\
\end{longtable}
\begin{tipbox}[提示]
判断“程序本身有 bug”还是“只在调试器下才出错”时，就用 \key{Ctrl+F5} 跑一遍。调试堆会把已释放的内存填成 0xDD、未初始化的堆内存填成 0xCD，有些野指针在调试器下“幸运地”崩掉，脱离调试器反而活得很好，这类现象本身就很有价值。
\end{tipbox}

\subsection{调试符号（.pdb）与 Debug 配置}

.pdb（Program Database）保存符号名、类型信息以及源文件与行号的映射，它\textbf{不包含}源代码，也不包含可执行代码；编译器用 \code{/Zi} 生成它，链接器用 \code{/DEBUG} 把它与 PE 文件关联。调试器靠 .pdb 里的 GUID + Age 与二进制严格配对，所以“断点变成空心圆圈，提示‘当前不会命中断点，还没有为该文档加载任何符号’”，九成是这四种情况之一：工程没开 \code{/Zi} 或 .pdb 没被复制到 exe 旁边；exe 与 .pdb 来自两次不同的编译；断点所在模块尚未加载，或那段代码根本没被执行到；符号路径没配对。

用 \ui{调试 > 窗口 > 模块} 可以查看每个模块的符号加载状态，在模块上右键可以手动指定 .pdb。系统 DLL 的符号不必自己找：在 \ui{工具 > 选项 > 调试 > 符号} 里打开“Microsoft 符号服务器”（地址形如 \url{https://msdl.microsoft.com/download/symbols}）即可，第一次会慢一些，之后有本地缓存；在\textbf{即时窗口}里也能用 \code{.sympath}、\code{.reload} 命令调整路径并强制重载。

Visual Studio 默认的 Debug 配置对 C++ 施加了一组互相配合的开关，它们的组合决定了“调试器能看到多少真相”：

\begin{longtable}{L{2.0cm}L{5.4cm}L{7.2cm}}\toprule
\textbf{开关} & \textbf{含义} & \textbf{对调试的影响} \\ \midrule \endfirsthead
\toprule \textbf{开关} & \textbf{含义} & \textbf{对调试的影响} \\ \midrule \endhead \bottomrule \endlastfoot
\code{/Od} & 禁用优化 & 变量不会被优化掉，语句顺序与源码一致，单步不乱跳 \\
\code{/Zi} & 生成完整调试信息到 .pdb & 断点能绑定到行号，类型信息可展开 \\
\code{/ZI} & “用于编辑并继续”的调试信息 & 要边调试边改代码时必须用它 \\
\code{/RTC1} & 运行时检查（\code{/RTCs} + \code{/RTCu}） & 栈帧填充为 0xCC，未初始化局部变量一眼可见 \\
\code{/JMC} & 仅我的代码单步支持 & 默认开启，决定 \key{F11} 会不会进入库函数 \\
\code{/MDd} & 调试版 CRT & 启用调试堆与内存检查 \\
\end{longtable}
优化是调试最主要的敌人：内联让“逐语句”一下跨过好几个函数，寄存器分配让变量显示“无法获取值”，循环展开让行号来回跳。必须在 Release 下排查时，至少加上 \code{/Zi} 并保留链接器的 \code{/DEBUG:FULL}，再用 \ui{调试 > 窗口 > 反汇编} 辅助理解。几个“填充图案”值得记住：\textbf{0xCCCCCCCC} 是 Debug 下未初始化的栈内存（\code{/RTCs} 的功劳），\textbf{0xCDCDCDCD} 是未初始化的堆内存，\textbf{0xDDDDDDDD} 是已被释放的堆内存，\textbf{0xFDFDFDFD} 是调试堆块两端的守护字节。

\subsection{调试器窗口布局与常用快捷键}

\begin{lstlisting}[style=textstyle]
中断模式下的一屏（窗口之间可以自由停靠）
  +-- 左上：调用堆栈
  |     main()
  |     > BumpAll(int*, int)
  |     [外部代码]
  +-- 中间：源代码编辑器
  |     data[i] += 1;          <-- 黄色箭头（即将执行的语句）
  +-- 右上：监视 1 / 监视 2
  |     g_buf.canary
  |     found   0x0000..
  |     i       4
  +-- 下方选项卡：局部变量 / 自动 / 断点 / 异常设置 / 输出 / 即时窗口
  |               内存 1 / 反汇编
  +-- 顶部：调试工具栏 [继续] [全部中断] [停止] [重启]
        菜单：调试 > 窗口 > ...
\end{lstlisting}
黄色箭头表示“即将执行的语句”，它还没有被执行；中断模式下，这一屏里的窗口是你唯一的信息来源。

\begin{longtable}{L{2.6cm}L{4.4cm}L{2.6cm}L{4.6cm}}\toprule
\textbf{窗口} & \textbf{快捷键} & \textbf{窗口} & \textbf{快捷键} \\ \midrule \endfirsthead
\toprule \textbf{窗口} & \textbf{快捷键} & \textbf{窗口} & \textbf{快捷键} \\ \midrule \endhead \bottomrule \endlastfoot
断点 / 调用堆栈 & \key{Ctrl+Alt+B} / \key{Ctrl+Alt+C} & 异常设置 / 模块 & \key{Ctrl+Alt+E} / \key{Ctrl+Alt+U} \\
监视 1--4 & \key{Ctrl+Alt+W}，再按 1--4 & 局部变量 / 自动 & \key{Ctrl+Alt+V}，再按 \key{L} / \key{A} \\
即时 / 输出 & \key{Ctrl+Alt+I} / \key{Ctrl+Alt+O} & 线程 / 并行堆栈 & \key{Ctrl+Alt+H} / \key{Ctrl+Shift+D, S} \\
内存 / 寄存器 / 反汇编 & \key{Ctrl+Alt+M}（再按 1--4）/ \key{G} / \key{D} & 打开方式都在 \ui{调试 > 窗口} 下 & — \\
\end{longtable}
\section{断点体系}

\subsection{普通断点与断点图标}

\key{F9}（\ui{调试 > 切换断点}）在光标所在行插入或删除断点，也可以直接点击代码左侧的边距；\key{Ctrl+F9} 切换启用 / 禁用，\key{Ctrl+Shift+F9} 删除全部断点。图标状态本身就在传递信息：\textbf{红色实心圆}表示已启用且已绑定到代码；\textbf{空心圆}表示已禁用；\textbf{空心圆加警告}表示未绑定，悬停可看原因；\textbf{带加号的圆}表示还附带条件、命中次数、筛选器或标签。

\subsection{条件断点与命中次数}

右键断点选择“条件”（快捷键 \key{Alt+F9} 后按 \key{C}），或悬停断点时点齿轮图标，打开“断点设置”窗口，里面有“条件表达式 / 命中次数 / 筛选器”三项。条件表达式有两种语义：\textbf{为 true}，如 \code{i == 4}、\code{p != nullptr && p->id == 3}；\textbf{更改时}，只要表达式取值发生变化就中断——注意调试器不把第一次求值算作“变化”，所以它不会在第一轮循环就停下。命中次数则跳过前 N 次，可选“等于 / 大于 / 大于等于 / 是 N 的倍数”，查“第 1000 次迭代之后才出错”这类 bug 时它比条件表达式更快，因为不必每次迭代都求值。

\begin{warnbox}[注意]
条件表达式由调试器在目标进程里求值，每次命中都要算一遍，放在热循环里会显著拖慢程序。尽量只写纯变量比较，避免在条件里调用函数——带有副作用的表达式可能被调试器拒绝求值，或者干脆改变程序行为。
\end{warnbox}

\subsection{筛选器断点}

同一个函数被主线程和 8 个工作线程反复调用时，条件表达式帮不上忙，该筛选器上场了。它支持 \code{ThreadId}、\code{ThreadName}、\code{ProcessId}、\code{ProcessName}、\code{MachineName} 五个字段，例如 \code{ThreadName = "Worker"}、\code{ProcessId = 6789}；多个子句用 \code{&}（与）、\code{||}（或）、\code{!}（非）和括号组合，字符串值必须加双引号。这是多线程调试里性价比最高的一招：先给线程起个名（见 5.3），再用 \code{ThreadName} 把断点限定住。

\subsection{函数断点}

\ui{调试 > 新建断点 > 函数断点}（\key{Ctrl+K, B}）适合“知道函数名但不想翻文件”或“同名重载函数想全部拦下”的场景。函数名可以写成全限定名 \code{MyNamespace::MyClass::Method} 来消除歧义，可以带参数类型 \code{Method(int, char*)} 来区分重载，可以用 \code{util.dll!Method} 或 \code{{Method, , util.dll}} 指定模块，也可以用 \code{{Method, , util.dll}+2} 跳过函数序言直接停在函数体，甚至写成 \code{((MyClass*)0x0000000123456780)->Method} 只在这个对象实例上中断。调用堆栈窗口里也可以右键某一帧的函数名插入断点——“我只想在这个调用路径上停下来”时非常顺手。

\subsection{数据断点（监视地址被改写）}

数据断点监视的是\textbf{某块内存被写入}这件事本身，而不是某一行代码；当变量地址在运行时才确定，或者你怀疑有人偷偷改了一块内存时，它是唯一的选择。创建方式有两种：\ui{调试 > 新建断点 > 数据断点}，在“地址”里填求值结果为地址的表达式（如 \code{&g_buf.canary}）；或者在自动、监视、局部变量窗口里右键变量选择“值更改时中断”。之后在“字节数”下拉框里选择要监视的宽度。限制条件必须记牢：

\begin{longtable}{L{3.0cm}L{12.0cm}}\toprule
\textbf{限制} & \textbf{说明} \\ \midrule \endfirsthead
\toprule \textbf{限制} & \textbf{说明} \\ \midrule \endhead \bottomrule \endlastfoot
只在本机 C++ 生效 & 原生代码专属；托管代码的数据断点限制更多，静态变量、结构体字段、带 DebuggerTypeProxy 特性的类都不支持 \\
只在写入时中断 & \textbf{读取不中断}；排查数组越界读必须换手段（条件断点、ASan 之类） \\
数量上限 & x64 / x86 最多 4 个，ARM64 最多 2 个，ARM 最多 1 个（由硬件调试寄存器数量决定） \\
宽度与对齐 & x86 下最宽 4 字节，x64 下最宽 8 字节；地址需按监视宽度对齐——监视 4 字节要求 4 字节对齐，监视 8 字节要求 8 字节对齐 \\
无效场景 & 别的进程改这块内存、内存被内核态改写（例如 \code{ReadFile} 填缓冲区）、共享内存，都不会触发 \\
\end{longtable}
\begin{warnbox}[注意]
数据断点绑定的是\textbf{绝对地址}，而每次重新调试变量地址几乎都会变，所以这类断点在调试会话结束时会被自动禁用，下次必须重新设置。同理，设在局部变量上的数据断点在该函数返回后地址就失效了，应该在返回前删除或禁用它。
\end{warnbox}

\subsection{临时断点、依赖断点与跟踪点}

\begin{itemize}
  \item \textbf{临时断点}：在边距右键选择“插入临时断点”，或在断点设置里勾选“命中一次后删除”，它只拦一次，等价于“把 \key{F5} 按到那一次为止”的自动化版本。
  \item \textbf{依赖断点}：在断点设置里勾选“仅当以下断点命中时启用”并选择一个前置断点，解决“这个工具函数被调用几千次，我只关心从某个入口进来的那一次”，比在条件表达式里写复杂的上下文判断干净得多。
  \item \textbf{跟踪点}（tracepoint）：断点不一定非要停下。设置“操作”让它向\textbf{输出}窗口打印一条带表达式的消息而程序继续跑，相当于一个随时可开关、不污染代码的日志语句。
\end{itemize}

\subsection{断点标签、导入导出与断点窗口}

\ui{调试 > 窗口 > 断点} 是所有断点的管理台，可以按名称、条件、命中次数、语言、标签排序和搜索，支持批量启用、禁用、删除。右键断点选择“编辑标签”，一个断点可以有多个标签：把和越界相关的断点都打上 \code{oob}，排查结束后按标签筛选、一次性清掉。断点窗口工具栏还提供“导出断点”与“导入断点”（XML 文件），把一组精心设计的断点随代码提交到仓库是团队协作里非常实用的技巧。另外记住：断点条件失效时（语法错误或语义不成立）调试器\textbf{仍然会中断}，只是给出警告——别指望它“静默跳过”。

\section{单步执行}

\subsection{单步命令与快捷键}

\begin{longtable}{L{3.6cm}L{4.4cm}L{6.6cm}}\toprule
\textbf{快捷键} & \textbf{命令} & \textbf{行为} \\ \midrule \endfirsthead
\toprule \textbf{快捷键} & \textbf{命令} & \textbf{行为} \\ \midrule \endhead \bottomrule \endlastfoot
\key{F10} & 逐过程（\ui{调试 > 逐过程}） & 执行当前行，不进入被调函数 \\
\key{F11} & 逐语句（\ui{调试 > 逐语句}） & 进入被调函数；无源码时进入反汇编 \\
\key{Shift+F11} & 跳出（\ui{调试 > 跳出}） & 执行到当前函数返回 \\
\key{Ctrl+F10} & 运行到光标处 & 在当前行下临时断点并启动调试 \\
\key{Ctrl+Shift+F10} & 设置下一条语句 & 改变执行点，不执行中间代码 \\
\key{F5} / \key{Shift+F5} & 继续 / 停止调试 & 运行到下一个断点 / 结束会话 \\
\key{Ctrl+Shift+F5} & 重新启动 & 重启进程，断点保留 \\
\end{longtable}
\begin{tipbox}[提示]
这一对快捷键常被记反：\textbf{运行到光标处是 \key{Ctrl+F10}}，设置下一条语句才是 \key{Ctrl+Shift+F10}。在编辑状态下右键代码行选择“运行到光标处”会同时完成“下临时断点 + 启动调试”两件事，比“先 \key{F9} 再 \key{F5}”更快；中断模式下把鼠标悬停在某行，行左侧会出现一个绿色的“运行到单击处”按钮，点它等价于运行到该行，这个命令在调用堆栈窗口里同样可用。
\end{tipbox}

\subsection{设置下一条语句与“往回走”}

\ui{调试 > 设置下一条语句}把执行点（黄色箭头）移到光标所在行，也可以直接用鼠标拖动黄色箭头。两种典型用途：跳过一次已知会失败的调用，看后面的逻辑是否正常；重放刚才那段代码，反复观察同一次执行的细节。但要清楚它的代价——它只改变指令指针，不撤销任何副作用：变量不会还原，已分配的内存不会回收，已析构的对象不会复活。对 C++ 而言，“往回跳”越过构造与析构、RAII 作用域、\code{std::unique_ptr} 的所有权转移，极易制造假象；它适合观察纯计算代码，不适合跨越资源管理边界。

\subsection{“仅我的代码”对单步的影响}

“仅我的代码”（Just My Code）在 \ui{工具 > 选项 > 调试 > 常规} 里开关，默认开启，它对 C++ 的影响是：

\begin{itemize}
  \item \textbf{调用堆栈}：非用户代码被折叠成一个灰色的 \code{[外部代码]} 帧，右键选择“显示外部代码”可以展开；
  \item \textbf{单步}：在非用户代码上按 \key{F11} 会被“跨过”，直接跑到下一行用户代码，\key{Shift+F11} 则直接跑到当前帧之外的下一行用户代码；
  \item \textbf{判定依据}：对调用堆栈而言，符号里没有源码信息、或源码文件不存在的函数算非用户代码；对单步而言判定更严格——对应 .pdb 没有加载的函数就算非用户代码。所以“某个第三方库明明有源码却跨不过去”，通常就是它的 .pdb 没加载，去 \ui{调试 > 窗口 > 模块} 确认一下；
  \item 单步相关的编译开关是 \code{/JMC}，C++ 工程默认开启；用较老的编译器产出的二进制只能靠 \code{.natstepfilter} 文件定制。
\end{itemize}

想进入 STL 内部看 \code{std::vector::push_back} 到底做了什么，就得关掉这个选项。调试器还支持用 \code{.natjmc} / \code{.natstepfilter} 两个 XML 文件精细指定“哪些模块、源文件、函数算外部代码”“哪些函数一律跨过”：可以放在 Visual Studio 安装目录的 \codeP{Common7\textbackslash{}Packages\textbackslash{}Debugger\textbackslash{}Visualizers} 下，也可以放在用户文档目录的 \code{Visualizers} 文件夹里，较新版本还允许直接放进项目或解决方案。

\begin{warnbox}[注意]
C++ 下即使开启“仅我的代码”，\textbf{异常仍然会在抛出的地方停下}，无论在用户代码还是外部代码里；“异常设置”里的“用户未处理”选项对 C++ 不适用。这一条经常让从 C\# 转过来的同学困惑。
\end{warnbox}

\section{查看数据}

\subsection{数据提示、自动窗口与局部变量}

中断模式下把鼠标悬停在变量上会弹出数据提示（DataTip），可以展开对象、查看 STL 容器的元素；点图钉可以把它“固定”在编辑器上，之后单步时会跟着更新，右键菜单里还能“添加监视”“编辑值”“复制表达式”。数据提示里的值同样支持格式说明符（见 4.3）。

\begin{longtable}{L{2.4cm}L{4.6cm}L{7.6cm}}\toprule
\textbf{窗口} & \textbf{打开方式} & \textbf{内容} \\ \midrule \endfirsthead
\toprule \textbf{窗口} & \textbf{打开方式} & \textbf{内容} \\ \midrule \endhead \bottomrule \endlastfoot
自动 & \ui{调试 > 窗口 > 自动} & 当前行及之前若干行用到的变量；C++ 下是\textbf{前 3 行} \\
局部变量 & \ui{调试 > 窗口 > 局部变量} & 当前作用域内的全部局部变量，含 \code{this} \\
监视 1--4 & \ui{调试 > 窗口 > 监视 > 监视 1} 等 & 手动输入的表达式，跨作用域保留 \\
\end{longtable}
自动窗口的内容随执行点滚动，适合“顺着代码看”；局部变量是当前作用域的全景；监视窗口是唯一能长期盯住某个表达式的窗口——变量离开作用域时它不会消失，只是变灰。

\subsection{监视窗口与即时窗口}

监视窗口接受的不只是变量名，任何调试器能求值的表达式都可以：

\begin{lstlisting}[style=textstyle]
p / *p / p->next->next     指针本身、解引用、链式追踪
&g_buf.canary              取地址，配合数据断点
(char*)p / sizeof(...)     强制转换、类型大小
arr,[4]                    把数组或指针按 4 个元素展开
\end{lstlisting}
\ui{调试 > 窗口 > 即时窗口}（\key{Ctrl+Alt+I}）更进一步：它除了求值，还能\textbf{执行}语句和函数调用，例如 \code{BumpAll(g_buf.data, 4)}、\codeP{printf("%d\textbackslash{}n", i)}、\code{SetLastError(0)}。这让它成为验证假设的利器——与其改代码重编译，不如在即时窗口里直接调用一次可疑函数；代价是它会真的改变进程状态，调用有副作用的函数容易把现场搞乱。

\begin{keybox}[实战场景]
怀疑某个对象的成员在某个时刻已经被损坏，就在监视窗口里写 \code{&obj} 记住地址，再写 \code{obj} 展开看内容；想看它何时被改写，直接右键该表达式选“值更改时中断”，比在源码里到处下断点快得多。
\end{keybox}

\subsection{格式说明符与并行监视}

在监视、自动、局部变量窗口里，变量名后面加一个逗号和说明符就能改变显示方式；监视窗口还有个技巧：\textbf{输入逗号后会弹出可用的说明符列表}，记不住时直接选。

\begin{longtable}{L{3.2cm}L{6.8cm}L{4.6cm}}\toprule
\textbf{说明符} & \textbf{效果} & \textbf{示例} \\ \midrule \endfirsthead
\toprule \textbf{说明符} & \textbf{效果} & \textbf{示例} \\ \midrule \endhead \bottomrule \endlastfoot
\code{x} / \code{X} / \code{xb} & 十六进制（带 / 不带 0x 前缀） & \code{flags,x} \\
\code{c} / \code{b} / \code{bb} / \code{o} & 字符 / 二进制 / 八进制 & \code{ch,c}、\code{mask,bb} \\
\code{e} / \code{f} / \code{g} & 科学计数 / 定点 / 自动浮点 & \code{ratio,f} \\
\code{s} / \code{s8} / \code{su} & 窄字符串 / UTF-8 / UTF-16 宽字符串 & \code{wstr,su} \\
\code{wc} / \code{wm} / \code{hr} & 窗口类标志 / Windows 消息号 / HRESULT 与 Win32 错误码 & \code{cls,wc}、\code{hres,hr} \\
\code{en} / \code{fe} / \code{!} & 枚举 / 位标志枚举 / 原始格式（忽略定制视图） & \code{state,en} \\
\code{,[N]} & 把指针当 N 个元素的数组 & \code{p,[10]}、\code{p,[0x20]}、\code{p,[count]} \\
\end{longtable}
两点提醒：数组长度现在推荐写 \code{,[N]} 形式，老资料里的 \code{p,10} 属于旧写法，长度也可以是表达式（如 \code{p,[count]}），调试动态缓冲区时非常方便；十六进制可以写成 \code{,x}，有些资料写作 \code{,hex}，具体以逗号下拉列表里列出的为准——官方文档还提示，一部分老式说明符（如 \code{d}、\code{h}、\code{u} 等）依赖旧的原生调试引擎，在新版调试器中不再生效。调试多线程或并行任务时，\ui{调试 > 窗口 > 并行监视 > 并行监视 1} 可以把同一个表达式按线程分行显示，一眼看出“哪个线程的值不对”，它常和并行堆栈配合使用。

\section{调用堆栈与线程}

\subsection{调用堆栈窗口}

\ui{调试 > 窗口 > 调用堆栈}（\key{Ctrl+Alt+C}）显示“我是怎么走到这里的”：最上面一行是当前函数，下面依次是调用者，最下面是 \code{main} 或线程入口。双击某一帧会切换调试上下文——局部变量、监视、\code{this} 都随之切换，但\textbf{程序并不会继续执行}。这一点是理解调用堆栈的关键：切换堆栈帧只是换了一个观察角度。右键菜单里的常用操作有插入断点、运行到光标处、显示外部代码、查看符号加载信息；有符号但找不到源码的帧会显示为“无可用源代码”。

\subsection{并行堆栈与线程窗口}

\ui{调试 > 窗口 > 并行堆栈}（\key{Ctrl+Shift+D, S}）把每个线程的调用栈并排画出，相同调用路径会被合并，于是“8 个线程都堵在同一个锁上”这种事一眼就能看出来；它可以在“线程”和“任务”两种视图间切换，也能显示或隐藏外部代码、按调用路径分组。\ui{调试 > 窗口 > 线程}（\key{Ctrl+Alt+H}）则列出进程内所有线程的 ID、类别、名称与当前位置，三个操作最实用：\textbf{切换线程}（双击一行，或用调试工具栏右侧的线程下拉框，切换后的调用堆栈与局部变量都属于该线程）、\textbf{重命名线程}（右键 > “重命名”，叫 \code{Worker-IO} 远比叫 \code{0x2A4C} 好记，而且线程名可以直接用在断点筛选器里）、\textbf{冻结 / 解冻线程}（暂时不让某个线程被调度，排查“到底哪个线程在破坏数据”时很有用）。

\subsection{多线程死锁的排查思路}

死锁的典型形态是 A 等 B 的锁、B 等 A 的锁，两个线程都“卡住但活着”。步骤是：程序卡住后先 \ui{调试 > 全部中断} 打断所有线程；打开 \ui{调试 > 窗口 > 并行堆栈} 看有几个线程停在锁操作上；逐个线程查看调用堆栈，分别记录“它正在等待的对象”和“它已经持有的对象”；然后找环——在监视窗口里展开 \code{CRITICAL_SECTION} 的 \code{OwningThread}、\code{LockCount}、\code{LockSemaphore} 等字段确认谁在持有（\code{std::mutex} 的内部实现随版本而变，不好直接读，这也是“自己包一层互斥量并记录持有者”这种做法的现实理由）。

\begin{lstlisting}[style=textstyle]
线程 A:  Lock(mutexA)  <- 成功  -->  Lock(mutexB)  <- 阻塞，等 B 释放
线程 B:  Lock(mutexB)  <- 成功  -->  Lock(mutexA)  <- 阻塞，等 A 释放
             ^
             +-- 互相等待，谁都不放手  =>  死锁
\end{lstlisting}
\begin{tipbox}[提示]
预防比排查便宜。约定全局统一的加锁顺序、用 \code{std::scoped_lock} 一次锁住多个互斥量（它内部带死锁避免算法）、避免在持锁期间调用可能回调用户代码的函数，这三条能消灭绝大多数死锁。
\end{tipbox}

\section{内存与底层调试}

\subsection{内存窗口}

\ui{调试 > 窗口 > 内存 > 内存 1}（\key{Ctrl+Alt+M}，\key{1}）直接查看进程内存，地址栏可以输入 \code{&g_buf}、\code{p}、\code{0x0000000123456780} 这类表达式，右键或工具栏可以切换每列的字节数（1 / 2 / 4 / 8）与“自动重新计算”。查指针与数组时：想知道 \code{p} 指向什么，就把 \code{p} 作为地址填进去，再按 1 字节或 4 字节解释；想确认越界有没有踩坏相邻变量，直接看 \code{&g_buf.canary} 附近的内存，并对照填充图案（0xCD / 0xDD / 0xFD）判断是“未初始化”还是“已被释放”；单纯想看字符串，在监视窗口里用 \code{,s} / \code{,su} 比在内存窗口里数字符省事得多。

\subsection{反汇编窗口与寄存器窗口}

\ui{调试 > 窗口 > 反汇编}（\key{Ctrl+Alt+D}）显示当前堆栈帧对应的汇编，并尽量把源码行交错显示；它能做到源码做不到的事：在优化过的 Release 代码里单步、看清编译器是否把某个分支优化掉了、确认崩溃时到底是哪条指令出了问题，右键菜单可以在“显示字节代码”“显示符号名”“显示行号”“转到源代码”之间切换。\ui{调试 > 窗口 > 寄存器}（\key{Ctrl+Alt+G}）显示通用寄存器与标志位，值可以直接编辑，也可以在监视窗口里用 \code{@rax}、\code{@rip} 这样的形式读取。x64 调用约定值得顺手记住：前 4 个整型或指针参数放在 RCX、RDX、R8、R9，浮点参数放在 XMM0--XMM3，返回值在 RAX——在反汇编里看到 \code{mov rax, [rcx]}，就基本能猜到是在解引用第一个参数。

\begin{keybox}[实战场景]
一个只在 Release 下出现的崩溃，断点绑不上。做法是：在异常设置里勾上对应的 Win32 异常，让它在崩溃瞬间中断，然后打开反汇编窗口看当前指令，再用寄存器窗口里的地址配合内存窗口确认指针是否为空、是否已被释放。
\end{keybox}

\section{异常处理}

\subsection{异常设置窗口与“两类异常”}

\ui{调试 > 窗口 > 异常设置}（\key{Ctrl+Alt+E}）决定“哪些异常会让调试器停下来”，原生 C++ 调试下主要看两个节点：\textbf{Win32 异常}，即操作系统层面的异常代码，常见的有 0xC0000005（访问冲突）、0xC0000094（整数除零）、0xC00000FD（栈溢出）、0xC0000374（堆损坏）；\textbf{C++ 异常}，即 C++ 的 \code{throw}，可以按异常类型添加过滤条件。

异常调试最容易混淆的是“首次异常”与“用户未处理异常”：

\begin{longtable}{L{3.4cm}L{3.6cm}L{3.4cm}L{3.8cm}}\toprule
\textbf{概念} & \textbf{时机} & \textbf{此时的堆栈} & \textbf{用途} \\ \midrule \endfirsthead
\toprule \textbf{概念} & \textbf{时机} & \textbf{此时的堆栈} & \textbf{用途} \\ \midrule \endhead \bottomrule \endlastfoot
首次异常（第一轮） & 异常刚被抛出，处理程序还没运行 & 完整，现场最新鲜 & 找“到底是谁抛的” \\
用户未处理异常（第二轮） & 没有任何处理程序接管，即将终止进程 & 可能已被栈展开破坏 & 找“为什么没人接住” \\
\end{longtable}
对 C++ 来说，栈展开（unwinding）会调用沿途对象的析构函数，等你在“未处理异常”处中断时，局部对象可能已经析构，变量的值已经不可信，所以\textbf{排查 C++ 异常要优先在首次异常处中断}。

\begin{warnbox}[注意]
默认情况下“仅我的代码”会把外部代码抛出的异常隐藏起来，只在输出窗口留一行消息。要让调试器在系统库或第三方库里抛出异常时也停下，需要关闭“仅我的代码”，并在异常设置里勾上对应类型。
\end{warnbox}

\subsection{C++ 异常、SEH 与“在 throw 处中断”}

C++ 异常（\code{throw std::runtime_error(...)}）是语言层面的，由运行时函数负责抛出与展开；SEH（\code{__try} / \code{__except} / \code{__finally}）是操作系统层面的，访问冲突、除零这类硬件异常天然就是 SEH 异常，\code{__except} 的过滤表达式可以用 \code{GetExceptionCode()} 取到异常代码。当两者混用时，\code{/EHsc} 与 \code{/EHa} 的行为差别会成为 bug 的来源：\code{/EHa} 让 \code{catch(...)} 能接住 SEH 异常，但也可能掩盖真正的访问冲突——调试时如果发现“该崩的地方没崩”，先检查这个开关。

想让调试器在 \code{throw} 处中断，有三种做法，精确度递增：在异常设置窗口的 \textbf{C++ 异常} 节点上勾选“抛出时”中断（覆盖面最广，但库内部的正常异常流也会被拦下，比较吵）；在具体的 \code{throw} 语句所在行下断点（精确，但只覆盖你找得到的那几处）；\textbf{在函数断点里输入 \code{_CxxThrowException}}——MSVC 的所有 C++ \code{throw} 最终都会走到这个运行时函数，因此它等价于“任何地方 throw 都中断”，还不会漏掉模板和库代码里的抛出，配合调用堆栈就能立刻看出是谁抛的、抛之前发生了什么。

\section{编辑并继续 / 热重载}

\subsection{前提条件}

“编辑并继续”（Edit and Continue）现在统一叫\textbf{热重载}（Hot Reload），但对 C++ 而言底层能力仍是编辑并继续。它只在\textbf{调试版本}下工作，且需要：编译用 \code{/ZI}（“用于编辑并继续”的调试信息格式），链接器启用增量链接 \code{/INCREMENTAL}；不能用 \code{/Zo}（增强优化调试）编译，那会\textbf{关闭}热重载；通常要求从 Visual Studio 用 \key{F5} 启动（附加到进程的场景一般不支持）。相关选项位于 \ui{工具 > 选项 > 调试 > .NET/C++ 热重载}：“启用热重载”“继续时自动应用更改（仅限本机）”“对陈旧代码发出警告”“不调试时启用”“保存文件时应用更改”“日志详细程度”，其中“继续时自动应用更改”默认开着。手动应用用 \ui{调试 > 应用代码更改}，它只应用修改、\textbf{不继续执行}；应用过程中想反悔，用 \ui{调试 > 停止应用代码更改}（该菜单项只在应用过程中可见），此时所有更改都不会提交。

\subsection{能改什么，不能改什么}

\begin{longtable}{L{4.6cm}L{2.4cm}L{7.6cm}}\toprule
\textbf{修改类型} & \textbf{能否热应用} & \textbf{说明} \\ \midrule \endfirsthead
\toprule \textbf{修改类型} & \textbf{能否热应用} & \textbf{说明} \\ \midrule \endhead \bottomrule \endlastfoot
函数体内的语句、局部变量初值与表达式 & 可以 & 最常用的场景；执行点可能被移动，C++ 会弹窗提示 \\
新增局部变量、修改当前函数的调用者或被调用者 & 有限制 / 暂缓 & 给\textbf{在调用堆栈上}的函数新增超过 64 字节变量时无法立即应用，会退化为“陈旧代码”；改了调用者则先执行旧代码 \\
修改类、结构体的成员与布局、函数签名、增删成员函数或全局变量 & 不可以 & 会让已有对象的内存布局失效或需要重新编译链接，必须重启 \\
改宏、改头文件里的类型定义、换库或改构建选项 & 不可以 & 属于构建配置变更 \\
\end{longtable}
无法应用时，Visual Studio 会明确提示“此更改需要重新生成 / 重新启动调试”。看到这句话就不要再试图绕过，老老实实 \key{Ctrl+Shift+F5} 重启。

\subsection{陈旧代码与执行点漂移}

两个概念值得单独记住：\textbf{陈旧代码}（stale code）指修改无法立即应用时（编辑了当前函数的调用者，或者给栈上的函数加了太多新变量），调试器继续执行旧代码，并把它显示在一个标题类似 \code{enc25.tmp} 的临时源文件窗口里，而你编辑后的文件仍在原窗口——此时不要编辑那个临时窗口；\textbf{执行点漂移}指某些修改会让黄色箭头的位置发生变化，C++ 下调试器会弹窗提示，需要你确认位置是否正确，不对就用“设置下一条语句”手动纠正。

\begin{tipbox}[提示]
热重载适合改打印、改分支条件、修正一两个表达式这类“微调”。如果一次要动十几行逻辑或者改数据结构，直接重启调试往往比和热重载较劲更快。
\end{tipbox}

\section{高级调试手段}

\subsection{附加到进程与远程调试}

\ui{调试 > 附加到进程}（\key{Ctrl+Alt+P}）适用于程序已经在跑、复现步骤又很长的场景：先让程序进入“即将出错”的状态，再附加，然后手动触发问题。“附加到”一栏要选对类型——纯原生程序选“本机代码”，混了托管代码选“自动”或“混合”，选错的典型症状是断点变成空心圆、调用堆栈里全是地址；64 位的 Visual Studio 可以调试 32 位和 64 位进程，但反向不行。

目标机器不能装 Visual Studio 时用远程调试：先把与目标机架构匹配的远程调试器（\code{msvsmon.exe}，位于 Visual Studio 安装目录的 \codeP{Common7\textbackslash{}IDE\textbackslash{}Remote Debugger\textbackslash{}<arch>} 下）复制到目标机并运行，记下它监听的端口；然后在开发机上用 \ui{调试 > 附加到进程}，连接类型选“远程（Windows 身份验证）”或“远程（无身份验证）”，目标填 \code{机器名:端口}；也可以在 \ui{项目属性 > 调试} 里把“要启动的调试器”设为“远程 Windows 调试器”、“远程服务器”填目标机名，之后直接 \key{F5} 就是远程启动。要点是：目标机上必须有与二进制\textbf{完全匹配}的 .pdb，源码路径也要能被开发机找到（否则只能看反汇编），防火墙需要放行 msvsmon 的端口。

\subsection{转储文件调试}

崩溃发生在别人的机器上、现场无法复现时，转储文件（.dmp）是唯一的事后证据：

\begin{lstlisting}[style=shellstyle]
# 让 procdump 在进程崩溃(未处理异常)时抓一份完整转储, 也可用任务管理器右键"创建转储文件"
.\procdump.exe -accepteula -ma -e -x .\dumps .\demo.exe
# 检查刚才那次运行的退出代码, 0xC0000005 表示访问冲突
$p = Start-Process .\demo.exe -PassThru -Wait; "退出代码: {0}" -f $p.ExitCode
\end{lstlisting}
打开转储文件用 \ui{文件 > 打开 > 文件}，选中的 .dmp 会以调试会话的形式打开，调试器直接停在异常发生的那一帧；之后的操作和在线调试几乎一样，调用堆栈、线程窗口、局部变量、并行堆栈都能用。两个前提：符号路径要配好（至少指向自己程序的 .pdb），并且转储文件必须“完整”到包含你需要的线程栈与堆内存——小型转储体积小，但可能看不到局部变量的值。

\subsection{OutputDebugString、TRACE 与 assert}

在调试器之外，最轻量的观察手段是让程序自己说话：

\begin{lstlisting}[style=cppstyle]
static void DebugLog(const char* fmt, ...) {        // 记得用 #ifdef _DEBUG 包住
    char buf[512]; va_list a; va_start(a, fmt); _vsnprintf_s(buf, _TRUNCATE, fmt, a); va_end(a);
    ::OutputDebugStringA(buf);                      // <windows.h>; 输出到 "输出" 窗口的 "调试" 页
}
_RPT1(_CRT_WARN, "count = %d\n", count);            // CRT 宏, 自动带文件名与行号
TRACE("value = %d\n", value);                       // MFC / ATL 宏, 仅在 _DEBUG 下有效
assert(head != nullptr && "head 不应为空");          // 断言表达式里绝不要写有副作用的代码
_ASSERTE(index >= 0 && index < count);              // _ASSERT 只输出信息, _ASSERTE 输出表达式文本
\end{lstlisting}
输出出现在 \ui{调试 > 窗口 > 输出} 的“调试”页，脱离调试器时可以用 Sysinternals 的 DebugView 捕获；这类调用本身有开销，一定要用 \code{#ifdef _DEBUG} 包住，或者改用 ETW。断言则是“开发期的合同检查”：

Debug 下断言失败会弹出断言对话框（中止 / 重试 / 忽略），点“重试”会\textbf{直接进入调试器}并停在失败的那一行，这是从断言跳进调试会话的标准入口；定义 \code{NDEBUG} 后 \code{assert} 会被完全消除，所以 \code{assert(++i < n)} 是经典错误。断言对话框的行为可以用 \code{_CrtSetReportMode} / \code{_CrtSetReportHook} 改写，把消息重定向到日志文件，适合自动化测试环境；编译期的检查则交给 \code{static_assert}。

\subsection{调试时的环境变量与命令行参数}

\ui{项目属性 > 调试} 这一页决定了 \key{F5} 到底怎样启动你的程序，却经常被忽略：

\begin{longtable}{L{4.0cm}L{5.4cm}L{5.2cm}}\toprule
\textbf{属性} & \textbf{作用} & \textbf{注意事项} \\ \midrule \endfirsthead
\toprule \textbf{属性} & \textbf{作用} & \textbf{注意事项} \\ \midrule \endhead \bottomrule \endlastfoot
命令 / 命令参数 & 要启动的可执行文件，以及传给 \code{main(argc, argv)} 的参数 & 默认 \code{$(TargetPath)}；参数含空格要加引号，改动后必须重启调试才生效 \\
工作目录 & 进程的当前目录 & 影响相对路径读写；与 exe 所在目录不同会导致 DLL 找不到 \\
环境 / 合并环境 & 追加或覆盖环境变量，每行一个 \code{名称=值}；是否与系统环境合并 & 常用于给 \code{PATH} 追加 DLL 目录；合并设为“否”时环境更干净，便于排除干扰 \\
调试器类型 / 附加 / 远程服务器 & 自动、本机、混合、仅托管；是否附加到已有进程；远程调试的机器名与端口 & 纯 C++ 选“本机”；“附加”为“是”时 \key{F5} 不启动新进程；远程调试见 9.1 \\
\end{longtable}
环境变量是排查“本机能跑、换台机器就少 DLL”的第一手段。与其把 DLL 拷来拷去，不如在“环境”里加一行 \codeP{PATH=%PATH%;D:\textbackslash{}libs\textbackslash{}bin;\$(OutDir)}。CMake 工程不用这套属性页，而是在 \code{launch.vs.json} 里描述同样的东西。命令行构建与调试符号的关系也很直接：

\begin{lstlisting}[style=shellstyle]
rem 在 "x64 Native Tools Command Prompt for VS" 里执行: 带符号的调试版本
cl /nologo /Zi /Od /EHsc /W4 /MDd main.cpp
\end{lstlisting}
\section{实战示例}

下面这段代码故意埋了两个经典 bug：\textbf{数组越界写}与\textbf{空指针解引用}。把它存成 \code{main.cpp}，跟着做一遍，就能把前面几节的内容串起来。

\begin{lstlisting}[style=cppstyle]
// cl /nologo /Zi /Od /EHsc /W4 /MDd main.cpp
#include <cstdio>
struct Buffer { int data[4]; int canary; };      // canary 在布局上紧跟 data[3] 之后
static Buffer g_buf = { { 10, 20, 30, 40 }, 0x7FFFFFFF };
struct Node { int id; Node* next; };
static void BumpAll(int* data, int count) {
    for (int i = 0; i <= count; ++i) data[i] += 1;   // BUG 1: 应为 i < count, 越界踩到 canary
}
static Node* FindNode(Node* head, int id) {
    for (Node* p = head; p != nullptr; p = p->next) if (p->id == id) return p;
    return nullptr;                              // 找不到就返回空指针
}
int main() {
    printf("canary before = 0x%08X\n", g_buf.canary);
    BumpAll(g_buf.data, 4);
    printf("canary after  = 0x%08X\n", g_buf.canary);
    Node a{ 1, nullptr }, b{ 2, nullptr }; a.next = &b;
    Node* found = FindNode(&a, 3);               // 3 根本不存在
    printf("found id = %d\n", found->id);        // BUG 2: 空指针解引用
}
\end{lstlisting}
\subsection{用数据断点抓数组越界}

\code{canary} 的初值是 0x7FFFFFFF，如果运行后它变了，就说明有人越界写到了 \code{g_buf.data} 之外。步骤：

\begin{enumerate}
  \item 在 \code{BumpAll} 的循环体那行按 \key{F9} 下断点，\key{F5} 启动；第一次命中时在监视窗口加入 \code{i} 和 \code{&g_buf.canary}，记下地址（例如 \code{0x00007FF6A1B23014}）与当前 \code{i}；
  \item 继续按 \key{F5} 观察 \code{i} 递增到 3——\code{data[3]} 是最后一个合法元素，再往后就出界了；
  \item 删掉这个断点，改用\textbf{数据断点}：\ui{调试 > 新建断点 > 数据断点}，地址填 \code{&g_buf.canary}，字节数选 4；
  \item 按 \key{F5} 继续，程序会\textbf{恰好在越界写发生的那一瞬间}停下来，而不是等到 \code{printf} 打印出错误的 canary；
  \item 此时看调用堆栈：当前帧是 \code{BumpAll}，调用者是 \code{main}；监视窗口里 \code{i} 是 \textbf{4}，而 \code{count} 也是 4——\code{data[4]} 显然越界，\code{for (int i = 0; i <= count; ++i)} 里的 \code{<=} 就是元凶；
  \item 改成 \code{i < count} 后应用代码更改或重启调试，canary 应该稳定保持在 \code{0x7FFFFFFF}。
\end{enumerate}

如果没有数据断点，你只能靠“canary 变了”这个结果倒推，中间过程完全看不到；数据断点把这个过程压缩成了一次精确的中断。

\subsection{用断点 + 监视 + 调用堆栈抓空指针}

\begin{enumerate}
  \item 在 \codeP{printf("found id = %d\textbackslash{}n", found->id);} 这行按 \key{F9}，\key{F5} 运行到断点；
  \item 在监视窗口加入 \code{found}——值是 \code{0x0000000000000000}；展开 \code{found->id} 会提示无法读取内存，这就是崩溃的前兆；
  \item 打开 \ui{调试 > 窗口 > 调用堆栈}，双击 \code{main} 上面那一帧进入 \code{FindNode}；或者直接在 \code{return nullptr;} 那行下断点重新运行一次，看它是在什么输入下返回了空；
  \item 按 \key{F10} 逐过程执行到下一行（程序会在 \code{printf} 内部触发访问冲突）。这里可以做两件事：在 \ui{调试 > 窗口 > 异常设置} 里勾上 \code{Win32 Exceptions} 下的访问冲突（0xC0000005），让调试器在异常发生的第一时间中断；打开 \ui{调试 > 窗口 > 反汇编}，会看到类似 \code{mov eax, dword ptr [rax]} 的指令，而寄存器窗口里 RAX 是 0——机器层面的原因一目了然；
  \item 在即时窗口里验证假设：输入 \code{found != nullptr} 得到 \code{false}，再输入 \code{FindNode(&a, 3)} 确认它确实返回空；
  \item 修复思路有两种：在调用处判空后再解引用（最小改动），或者让 \code{FindNode} 返回引用、可选值之类的类型，把“可能找不到”编码进接口（更好的设计）。
\end{enumerate}

\begin{keybox}[实战场景]
真实项目里的崩溃往往比这复杂得多，但流程完全一样。记住这条“三窗口联动”的套路：\textbf{调用堆栈告诉你从哪来，监视窗口告诉你现在是什么状态，数据断点与异常设置告诉你是哪一刻变坏的}。三者配合，绝大多数内存类 bug 都能在一个调试会话里定位。
\end{keybox}

\section{小结}

\begin{itemize}
  \item \key{F5} 是“带调试器跑”，\key{Ctrl+F5} 是“纯跑一遍”；.pdb 靠 GUID + Age 与二进制配对，断点变成空心圆先查符号有没有加载，而 \code{/Zi} 与 \code{/Od} 是能在 Debug 下看清变量与单步的前提。
  \item 断点远不止“某一行”：条件、命中次数、筛选器、函数断点、数据断点、临时断点、依赖断点各有明确的适用场景；其中数据断点只在写入时中断、x64 下最宽 8 字节且需按宽度对齐、总数不超过 4 个，而断点窗口负责批量管理、打标签与导入导出。
  \item 单步用 \key{F10} / \key{F11} / \key{Shift+F11}，想跳位置用 \key{Ctrl+F10}（运行到光标处）和 \key{Ctrl+Shift+F10}（设置下一条语句，它只改指令指针、不撤销副作用）；看数据靠数据提示、自动与局部变量、监视 1--4、即时窗口，以及 \code{,[N]}、\code{,x}、\code{,su} 这类格式说明符。
  \item 并发问题用调用堆栈、并行堆栈与线程窗口这三件套，死锁的固定套路是“全部中断 + 逐线程看堆栈 + 找环”；只看机器层面时用内存、反汇编、寄存器窗口。
  \item 异常处理的关键是分清“首次异常”与“用户未处理异常”，对 C++ 在首次异常或 \code{_CxxThrowException} 处中断最有价值；热重载能改函数体、改不了数据布局；附加到进程、远程调试、转储文件、\code{OutputDebugString}、\code{assert} 与“项目属性 > 调试”构成了真实项目的完整工具箱。
\end{itemize}

\section{练习题}

\begin{enumerate}
  \item 写一个把 \code{int} 数组元素全部翻倍的函数，故意让循环多走一次；先用条件断点 \code{i == count} 找到越界，再改用数据断点监视数组之后的一个变量，比较两种方式各自的优缺点（提示：哪种能看见“写入的瞬间”？）。
  \item 用 \code{OutputDebugStringA} 加一个 \code{#ifdef _DEBUG} 包装宏，实现一个带文件名与行号的日志函数，让它的输出出现在\textbf{输出}窗口的“调试”页；再故意去掉 \code{_DEBUG}，确认 Release 下它完全不产生代码。
  \item 构造一个两线程死锁：线程 A 先锁 mutexA 再锁 mutexB，线程 B 反过来；程序卡住后用“全部中断”+ 并行堆栈定位，记录两个线程各自的调用堆栈，并说明如何用统一加锁顺序修复。
  \item 给第 10 节的例子加一个会抛出 \code{std::runtime_error} 的分支，分别用“异常设置里勾选 C++ 异常”和“函数断点 \code{_CxxThrowException}”两种方式让调试器在 throw 处中断，比较两种方式拦到的异常数量。
  \item 用 \ui{项目属性 > 调试} 的命令参数给程序传三个参数（其中一个含空格），并在“环境”里追加一个自定义变量，在代码中用 \code{getenv} 读出来；说明为什么改了参数之后必须重启调试。
  \item 在调试会话中修改一个函数体（比如改打印内容），用“应用代码更改”让它生效；然后试着给 \code{struct Buffer} 加一个新成员，观察 Visual Studio 的提示，并解释为什么这类修改必须重启。
\end{enumerate}

\section{相关笔记}

\begin{itemize}
  \item \textbf{本教程}：\textbf{Visual Studio 2026 使用教程}
  \item \textbf{上一章}：\textbf{代码编辑与 IntelliSense}
  \item \textbf{下一章}：\textbf{性能分析与诊断}
  \item \textbf{C++ 教程}：\textbf{Visual Studio 入门}
  \item \textbf{CMake}：\textbf{CMake 使用指南}
\end{itemize}
